% Generated draft: requires mathematical and source review.
\documentclass[11pt]{article}
\usepackage[margin=25mm]{geometry}
\usepackage{fontspec}
\setmainfont{Latin Modern Roman}
\usepackage{amsmath,amssymb,mathtools,unicode-math}
\setmathfont{Latin Modern Math}
\usepackage{enumitem,needspace,xurl,hyperref}
\setlength{\parindent}{0pt}
\setlength{\parskip}{5pt}
\setlength{\emergencystretch}{3em}
\newcommand{\N}{\mathbb{N}}
\newcommand{\Z}{\mathbb{Z}}
\newcommand{\Q}{\mathbb{Q}}
\newcommand{\R}{\mathbb{R}}
\newcommand{\C}{\mathbb{C}}
\newcommand{\F}{\mathbb{F}}
\newcommand{\A}{\mathbb{A}}
\newcommand{\PP}{\mathbb{P}}
\newcommand{\E}{\mathbb{E}}
\newcommand{\eps}{\varepsilon}
\newcommand{\dd}{\,\mathrm{d}}
\begin{document}
\section*{UnsolvedMath ID: 30003711}
\textbf{Source problem code:} catalogue entry
\subsection{Definitions and mathematical statement}
Let \(V=\C[x_0,x_1,x_2]_3\) be the vector space of homogeneous cubic polynomials, and set
\[
X=\{[F]\in\PP(V): C_F=\{p\in\mathbb{CP}^2:F(p)=0\}\text{ is smooth}\}.
\]
Here \([F]\) identifies nonzero scalar multiples of \(F\), and \(X\) carries the usual complex analytic topology. Thus the quotient specified in the statement is by \(\C^\times\); no quotient by projective coordinate changes is taken.

For \(p\in C_F\), let \(T_pC_F\) be its projective tangent line. A \emph{flex} is a point satisfying
\[
I_p(C_F,T_pC_F)\geq 3,
\]
where \(I_p\) denotes local intersection multiplicity. Define
\[
Y=\{([F],p)\in X\times\mathbb{CP}^2:p\text{ is a flex of }C_F\},\qquad
\pi:Y\longrightarrow X,\quad ([F],p)\longmapsto[F],
\]
with the subspace topology on \(Y\). The supplied statement identifies \(\pi\) as a degree-nine covering.

Use the unreduced convention for Schwarz genus:
\[
g(\pi)=\min\left\{k\geq1:\begin{array}{l}
X=U_1\cup\cdots\cup U_k,\quad U_i\subseteq X\text{ open},\\
\text{there exist continuous }s_i:U_i\to Y\text{ with }\pi\circ s_i=\operatorname{id}_{U_i}
\end{array}\right\},
\]
with value \(\infty\) if no finite such cover exists. Determine \(g(\pi)\).

The supplied literature assessment dated 2026-08-21 reports \(8\leq g(\pi)\leq9\) and labels the choice between these values open. This is supplied dated curation, not independently verified here. The record does not specify the genus normalization, so compatibility of those numerical bounds with the convention above requires review. Its background also uses the phrase ``moduli space''; the target here follows the explicitly stated scalar quotient, and applicability of that background assessment to this precise space requires verification.
\subsection{Short English statement}
Determine the minimum number of open sets needed to cover the parameter space of smooth complex plane cubics so that one can continuously choose a flex point on each set. Cubic equations are identified only up to nonzero scalar multiplication. The supplied dated assessment reports that the answer is either \(8\) or \(9\), subject to checking its space and genus convention.
\subsection{Sources}
\begin{itemize}
\item Local UnsolvedMath registry, supplied record with permanent ID \(30003711\), title ``Schwarz Genus of the Flex-Point Cover,'' source code \texttt{catalogue entry}, Oberwolfach Reports Open Problems collection. The supplied status is \texttt{open}; the supplied statement status is \texttt{corrected\_verified}, described as preserving mathematical content while repairing notation, typography, or omitted definitions. These are registry labels, not independent verification here. No individual proposer is supplied; the proposed year is 2018.
\item Supplied report citation: \emph{Topology of Arrangements and Representation Stability} (2019), Oberwolfach Reports, pp.~43--123, DOI \texttt{10.4171/owr/2018/2}. \url{https://doi.org/10.4171/owr/2018/2}. The supplied background attributes the problem collection to mathematical workshop reports and their open-problem discussions.
\item Supplied literature source: Chen--Wan, ``flex-point topological complexity.'' \url{https://arxiv.org/abs/2306.17303}.
\item Supplied literature source: Chen--Wan, \emph{Proceedings of the AMS}. \url{https://doi.org/10.1090/proc/17184}.
\item Literature assessment date: \texttt{2026-08-21}. The supplied assessment reports lower bound \(8\), upper bound \(9\), and an open choice between them. This is supplied dated curation derived from the OWR record, not independently verified; the record explicitly requests independent verification of time-sensitive status and cited claims.
\end{itemize}
\end{document}
